% sections/05_proof_sketches.tex
% JPHQ-04 — Proof sketches and proof obligations (OC/ETS-internal)

Proofs in this artifact are \emph{internal} to the
\OC/\kw{ETS.K\_EA.v1.1} formal system. Accordingly, we provide proof
sketches together with explicit \emph{proof obligations}, i.e., statements
that must be derivable from the assumed OC/ETS primitives, definitions,
and axioms. No external semantics, empirical assumptions, behavioral
models, or interpretive arguments are employed.

Each proof sketch follows a uniform logical structure:
\begin{enumerate}
  \item explicit structural assumptions stated entirely in OC/ETS terms;
  \item reduction to an incompatibility expressed solely via admissibility
  conditions on $\Om(\KEA)$, $\Ax(\KEA)$, $\Th(\KEA)$, $\Cyc(\KEA)$, or the
  phase-transition logic;
  \item contradiction with admissibility of states, cycles, or phase paths;
  \item conclusion: non-existence of admissible transformations, except via
  phases $\PHCO$ or $\PHST$.
\end{enumerate}

\subsection*{Proof Sketch for Theorem~\ref{thm:axis-incompatibility}}

\textbf{Assumptions.}
Any transformation of $\KEA$ necessarily requires the introduction of an
axis $A_{\mathrm{new}} \notin \Ax(\KEA)$ whose constraints are incompatible
with the existing axis set $\Ax(\KEA)$ under the threshold configuration
$\Th(\KEA)$.

\textbf{Reduction.}
Axis incompatibility implies that there exists no state satisfying
simultaneously the constraints induced by
$\Ax(\KEA)\cup\{A_{\mathrm{new}}\}$ under $\Th(\KEA)$.

\textbf{Contradiction.}
Admissibility of a transformation would require existence of at least one
state compatible with all axes and thresholds, i.e.,
$\Om(\KEA)\neq\varnothing$ under the extended axis set. This contradicts
the assumed incompatibility, which enforces
$\Om(\KEA)=\varnothing$.

\textbf{Conclusion.}
No admissible transformation exists; any realization of
$A_{\mathrm{new}}$ necessarily violates admissibility and is therefore
classified as $\PHCO$ or $\PHST$.

\textbf{Proof obligation.}
Derive within OC/ETS that incompatibility of axes under fixed thresholds
entails emptiness of the admissible state set.

\subsection*{Proof Sketch for Theorem~\ref{thm:threshold-barrier}}

\textbf{Assumptions.}
Any transformation requires a reconfiguration of thresholds $\Th(\KEA)$
such that no state satisfies all resulting constraints.

\textbf{Reduction.}
A threshold reconfiguration induces a modified admissibility condition on
$\Om(\KEA)$.

\textbf{Contradiction.}
Existence of an admissible transformed continuum requires
$\Om(\KEA)\neq\varnothing$ under the new threshold configuration. The
assumed reconfiguration enforces $\Om(\KEA)=\varnothing$, yielding a
direct contradiction.

\textbf{Conclusion.}
Transformation is formally impossible under the specified threshold
configuration.

\textbf{Proof obligation.}
Show within OC/ETS that satisfaction of threshold constraints is a
necessary condition for non-emptiness of $\Om(\KEA)$.

\subsection*{Proof Sketch for Theorem~\ref{thm:cycle-closure}}

\textbf{Assumptions.}
All admissible cycles $\Cyc(\KEA)$ become incompatible or undefined under
the postulated transformed structural type.

\textbf{Reduction.}
Absence of admissible cycles eliminates the closure required to maintain
continuity, enforcing $\kcont(\KEA)=0$.

\textbf{Contradiction.}
Any admissible transformed continuum must satisfy $\kcont(\KEA)>0$.
If no admissible cycle exists, continuity vanishes, contradicting
admissibility.

\textbf{Conclusion.}
Transformation necessarily results in phase $\PHST$ and is therefore
structurally impossible.

\textbf{Proof obligation.}
Establish within OC/ETS that existence of admissible cycles is required for
$\kcont(\KEA)>0$.

\subsection*{Proof Sketch for Theorem~\ref{thm:phase-incompatibility}}

\textbf{Assumptions.}
The OC/ETS phase-transition logic forbids any admissible path from the
current phase of $\KEA$ to a phase compatible with a transformed structural
type.

\textbf{Reduction.}
An admissible transformation presupposes existence of at least one allowed
phase-transition path consistent with the target structural type.

\textbf{Contradiction.}
If no such path exists, realization of a transformed structure would
violate the phase-transition constraints of OC/ETS.

\textbf{Conclusion.}
Transformation is impossible within the OC/ETS phase structure.

\textbf{Proof obligation.}
Demonstrate that admissible transformations require existence of at least
one allowed phase-transition path.

\subsection*{Proof Sketch for Theorem~\ref{thm:boundary-obstruction}}

\textbf{Assumptions.}
All trajectories required for transformation must cross $\dOm(\KEA)$ in a
manner not admitted by OC/ETS boundary conditions.

\textbf{Reduction.}
Boundary constraints restrict admissible trajectories in the state space
$\Om(\KEA)$.

\textbf{Contradiction.}
If every trajectory leading to a transformed structural type violates
boundary admissibility, no admissible path to such a structure exists.

\textbf{Conclusion.}
Transformation is structurally impossible.

\textbf{Proof obligation.}
Show within OC/ETS that trajectories crossing $\dOm(\KEA)$ outside
admissible boundary conditions are forbidden.

\begin{remark}
These proof sketches deliberately avoid completeness claims. They isolate
the minimal internal dependencies required for full formal proofs within
OC/ETS and thereby identify precisely where falsification of the No-Go
results would have to occur.
\end{remark}
